% s0: the theory in one sentence and five points, and how to read the boxes. ASCII only.
\begin{quote}\small
\textbf{The theory.} The mass gap of Yang--Mills theory is the energy cost of the cheapest way for the field to lean
toward changing its topology. It is carried by an identified shape at every scale, protected from the fine grain by
asymptotic freedom, and inherited all the way out to infinite space. If the theory is true, it says five things.
\textbf{(1)~The gap has a name:} in a small room it is carried by Helike, a twisting shape of charge~2 under the room's rotation symmetry, or by
her competitor, a breathing shape of charge~0, and as the room grows the carrier changes only at exact, countable
switches (\S4--5). \textbf{(2)~Topology decides who is lightest:} the energy landscape of the room is the squared slope
of the Chern--Simons function, and the competitor is exactly the direction that points toward the other winding vacuum,
which would explain why the scalar glueball is lightest (\S5; proved in the reduced model, while the full theory waits on one uncomputed term). \textbf{(3)~Tensor and scalar trade places at a
calculable point,} $s_0^*(t)=(t+2)(\sqrt{t+8}-\sqrt t)/(4\sqrt t)$, which is $1.1436$ at $t=3$ at harmonic order; the switches registered at stronger coupling rise from it (\S5). \textbf{(4)~The
fine grain cannot erase the note:} the infinitely many fast modes, at leading order, only rename the coupling, which is
asymptotic freedom ($b_0=22/3$ matched to $10^{-4}$) (\S6). \textbf{(5)~The infinite world is the last room:} infinite
space is the end of a family of spheres, and the gap is a section over that family that never touches zero (\S7--8).
Simulations and perturbation theory already show that the gap exists, that the scalar glueball is lightest and that
asymptotic freedom holds. What would be new is the reason behind all three, and that they are one fact seen from three
sides.
\end{quote}

\textbf{How to read this.} This is a narrative theory, not a proof, told as one story from the problem (\S1)
to what the gap is (\S10). Every important statement, proved or not, sits in a numbered box whose label gives its
standing: \textcolor{teal}{\textbf{PROVED}} (by me) and \textcolor{teal}{\textbf{PROVED (CITED)}} (by others),
\textcolor{sky}{\textbf{COMPUTED}} (numerically, not proved), \textcolor{gold}{\textbf{CONDITIONAL}} (proved from a
named hypothesis), \textcolor{plum}{\textbf{CONJECTURE}}, and \textcolor{muted}{\textbf{OPEN}}. Each box names the lines
of my register of statements~\cite{ClosedSet} it covers, and no box is stronger than its line. All 29 lines (29 is
merged into 23) appear below; \S10 tabulates where. The finer evidence, mine and
others', is in the companion blog post and the repository~\cite{HelikeReport}.
